\chapter{Game}\label{ch:K08_Game}

\section{Einführung in Pygame}
In diesem Kapitel entwickeln wir Schritt für Schritt ein eigenes 2D-Game mit \texttt{pygame-ce}. \texttt{pygame-ce} ist eine moderne, Community-gepflegte Variante von Pygame und wird genau gleich importiert mit \lstinline|import pygame as pg|. Sie eignet sich hervorragend, um in Python Grafiken, Animationen, Sound und Interaktionen umzusetzen.

Essentielle Bausteine eines Games sind:
\begin{itemize}
    \item Fenster (Auflösung, Titel) und Zeichenfläche (\enquote{screen})
    \item Game-Schleife mit \textbf{Ereignissen} (Tastatur, Maus) und \textbf{Zeitsteuerung} (\gls{FPS} / \enquote{refresh rate})
    \item Zeichnen: Hintergrund, Farben, geometrische Figuren, Bilder (\enquote{Sprites})
    \item \textbf{Zustände} und \textbf{Objekte} (z.B. Spieler als \lstinline|Rect|), \textbf{Bewegungen} und \textbf{Kollisionen}
    \item \textbf{Medien}: Bilder, Icon, Schrift, Sound / Musik
\end{itemize}

\begin{mydefinition}[Game-Schleife]
Ein Game läuft in einer Endlosschleife, bis das Programm beendet wird. In jeder Runde werden der Eingabestatus gelesen (\enquote{Events}), der interne Zustand aktualisiert (\enquote{Update}) und die Szene gezeichnet (\enquote{Render}). Eine \textbf{Clock} ($=$ Uhr) begrenzt die Bildwiederholrate (\gls{FPS}), sodass das Game stabil und gleichmässig läuft.
\end{mydefinition}

Um \tib{pygame-ce} in VS Code zu installieren, gehen Sie wie folgt vor:

\begin{enumerate}
    \item Öffnen Sie ein Terminal-Fenster in VS Code (\enquote{Terminal} $\rightarrow$ \enquote{New Terminal}).
    \item Installieren Sie pygame-ce mit pip:\\
          \lstinline|pip install pygame-ce|
    \item Überprüfen Sie die Installation, indem Sie eine neue Python-Datei \texttt{game\_test.py} mit folgendem Code ausführen:
          \begin{lstlisting}
        import pygame as pg
        print(pg.ver)
        \end{lstlisting}
          Dieser Code sollte, sofern \texttt{pygame-ce} korrekt installiert ist, die Versionsnummer von \texttt{pygame-ce} ausgeben. Damit ist \texttt{pygame-ce} einsatzbereit und Sie können mit den unten stehenden Übungen starten.
\end{enumerate}

\begin{myattention}
    Falls die Installation nicht funktioniert, erstellen Sie zuerst ein virtuelles Umfeld (venv) in VS Code und installieren Sie \texttt{pygame-ce} darin:
    \begin{itemize}
        \item Öffnen Sie ein Terminal-Fenster in VS Code (\enquote{Terminal} $\rightarrow$ \enquote{New Terminal}).
        \item Erstellen Sie ein virtuelles Environment:\\
              \lstinline|python3 -m venv venv|\\
              Aktivieren Sie das venv:\\
              \textbf{Windows:} \lstinline|venv\Scripts\activate|\\
              \textbf{macOS:} \lstinline|source venv/bin/activate|
        \item Führen Sie nun nochmals die obigen Schritte zur Installation von \texttt{pygame-ce} durch.
    \end{itemize}
\end{myattention}

\begin{myexample}[Minimale Game-Schleife]
Folgender Code zeigt den minimalen Aufbau eines Pygame-Programms mit Fenster, Game-Schleife und Ereignisverarbeitung. Kopieren Sie den Code in eine Datei \texttt{game\_intro.py} und führen Sie ihn in VS Code aus.
\lstinputlisting{Grundlagen_Info/00_Programmieren/Skript/Code/K08_Game/game_intro.py}
Ein schwarzes Fenster sollte erscheinen, das Sie mit dem Schliessen-Knopf beenden können.
\end{myexample}

\begin{myexercise}[Fenster-Titel setzen]
    Setzen Sie einen passenden Fenstertitel, indem Sie folgende Zeile abändern:

    \begin{lstlisting}
pg.display.set_caption("IHR TITEL HIER...")
\end{lstlisting}
\end{myexercise}

\begin{mychallenge}
    Laden Sie ein Fenster-Icon (\texttt{.png} oder \texttt{.jpg}):
    Laden Sie ein beliebiges Bild aus dem Internet herunter, speichern Sie es unter den Downloads und ziehen Sie es in Ihren Projektordner in VS Code (dort wo auch Ihre Python-Datei ist). Fügen Sie nun diese beiden Zeilen nach \lstinline|pg.display.set_mode(...)| ein:

    \begin{lstlisting}
icon = pg.image.load("path/to/icon.png")
pg.display.set_icon(icon)
\end{lstlisting}

    Wenn Ihr Ordner beispielsweise \texttt{Informatik/} heisst und das Bild unter
    \begin{center}
    \texttt{Informatik/Game/icon.png}
    \end{center}
    gespeichert ist, dann verwenden Sie
    \begin{center}
            \lstinline|pg.image.load("Game/icon.png")|.
    \end{center}
    Bei Windows ändert sich nur das Fenster-Icon, bei macOS wird das Icon nur im Dock angezeigt.
    \begin{myanswer}
        \lstinputlisting{Grundlagen_Info/00_Programmieren/Skript/Code/K08_Game/game_intro_window.py}
    \end{myanswer}
\end{mychallenge}

\begin{myexercise}[Hintergrund zeichnen]
    Füllen Sie den Hintergrund pro Frame mit einer Farbe mittels \lstinline|screen.fill((R,G,B))|. Experimentieren Sie mit zufällig generierten Farbtönen:

    \begin{lstlisting}
import random
# ...
# In der Game-Schleife, im Render-Abschnitt:
r = random.randint(0, 255)
g = random.randint(0, 255)
b = random.randint(0, 255)
screen.fill((r, g, b))
\end{lstlisting}
    Die Farbe muss in der Hauptschleife, aber noch vor \lstinline|pg.display.flip()|, gesetzt werden.

    \begin{itemize}
        \item Was ändert sich, wenn Sie die drei Zeilen für die Farbwerte \lstinline|r, g, b| vor die Hauptschleife setzen?
        \item Schreiben Sie den Code so um, dass die Farbe nur alle 10 Frames geändert wird.
    \end{itemize}


    \begin{myanswer}
        \lstinputlisting{Grundlagen_Info/00_Programmieren/Skript/Code/K08_Game/game_intro_background.py}
    \end{myanswer}
\end{myexercise}

\begin{myexercise}[Geometrische Figuren]
    Zeichnen Sie geometrische Formen: Rechteck, Kreis, Linie. Nutzen Sie dafür \lstinline|pg.draw.rect|, \lstinline|pg.draw.circle| und \lstinline|pg.draw.line|. Achten Sie darauf, \textbf{nach} dem Zeichnen \lstinline|pg.display.flip()| aufzurufen. Versuchen Sie, mittels folgender Befehle einen Apfel zu zeichnen:
    Verwenden Sie folgende Befehle:
    \begin{lstlisting}
pg.draw.rect(screen, rect_color, pg.Rect(rect_x, rect_y, rect_width, rect_height))
pg.draw.circle(screen, circle_color, (circle_x, circle_y), circle_radius)
pg.draw.line(screen, line_color, (start_x, start_y), (end_x, end_y), line_width)
\end{lstlisting}

    Dabei bezeichnen die Parameter Folgendes:
    \begin{itemize}
        \item \lstinline|screen|: die Zeichenfläche
        \item \lstinline|rect_color, circle_color, line_color|: Farbe als RGB-Tupel, z.\,B. \lstinline|(255, 0, 0)| für Rot
        \item \lstinline|pg.Rect(...)|: Rechteck-Objekt mit Position (x,y) und Grösse (width,height)
        \item \lstinline|(circle_x, circle_y)|: Mittelpunkt des Kreises
        \item \lstinline|circle_radius|: Radius des Kreises
        \item \lstinline|(start_x, start_y), (end_x, end_y)|: Start- und Endpunkt der Linie
        \item \lstinline|line_width|: Dicke der Linie in Pixeln (optional)
    \end{itemize}
\end{myexercise}

\begin{myexercise}[Sonnenuntergang erstellen]\label{myexercise:Sonnenuntergang}
    Erzeugen Sie eine einfache Szene mit Himmel, Sonne und Meer. Tipp: Zeichnen Sie einen Farbverlauf am Himmel, indem Sie in einer Schleife schmale horizontale Rechtecke in leicht unterschiedlichen Farbtönen zeichnen. Die Sonne ist ein Kreis; das Meer ein Rechteck in der unteren Hälfte. Das Meer sollte ebenfalls einen Farbverlauf haben (von hellblau zu dunkelblau). Verwenden Sie folgenden Code als Vorlage:
    \begin{lstlisting}
# Himmel (einfacher Gradient)
for y in range(0, HEIGHT//2):
	red = 100 + int(155 * (y / (HEIGHT//2)))  # 100 ... 255
	pg.draw.rect(screen, (red, 120, 180), (0, y, WIDTH, 1))

# Sonne
# IHR CODE HIER...

# Meer
# IHR CODE HIER...
\end{lstlisting}
    Optional: Lassen Sie die Sonne langsam sinken (Animation über mehrere Frames).
	\begin{myanswer}
		Diese Musterlösung finden Sie in \cref{sec:K09_sol}.
	\end{myanswer}
\end{myexercise}

\begin{myexercise}[Bewegungen (Tastatur)]\label{myexercise:Bewegungen}
    Erstellen Sie ein Spieler-\lstinline|Rect| (z.\,B. \lstinline|50 x 50|) und bewegen Sie es mit den Pfeiltasten. Nutzen Sie \lstinline|pg.key.get_pressed()| und begrenzen Sie die Bewegung auf das Fenster (\enquote{Screen Bounds}). Verwenden Sie folgende Code-Vorlage:
    \begin{lstlisting}
# Vor der Game-Schleife:
player = pg.Rect(100, 100, 50, 50) # Spieler-Rechteck erstellen
speed = 5 # Bewegungsgeschwindigkeit

# In der Game-Schleife:
keys = pg.key.get_pressed() # alle gedrückten Tasten abfragen
if keys[pg.K_LEFT]:
	player.x -= speed
if keys[pg.K_RIGHT]:
	player.x += speed
if keys[pg.K_UP]:
	player.y -= speed
if keys[pg.K_DOWN]:
	player.y += speed
player.clamp_ip(screen.get_rect())  # Rechteck innerhalb des Fensters halten

pg.draw.rect(screen, (0,0,0), player) # Spieler (Rechteck) zeichnen
\end{lstlisting}
	\begin{myanswer}
		Diese Musterlösung finden Sie in \cref{sec:K09_sol}.
	\end{myanswer}
\end{myexercise}

\begin{myexercise}[Kollisionen]\label{myexercise:Kollisionen}
    Legen Sie ein \lstinline|Rect| als \enquote{Ziel/Item} an (z.\,B. kleiner Kreis oder Block). Prüfen Sie eine Kollision mit \lstinline|player.colliderect(item)|. Bei Kollision: Position des Items neu zufällig setzen und optional Punkte zählen.
    \begin{lstlisting}
import random
# ...
# Vor der Game-Schleife:
item = pg.Rect(400, 300, 30, 30) # dieses Item soll gesammelt werden

# ...
# In der Game-Schleife, nach der Spieler-Bewegung:
if player.colliderect(item):
	item.topleft = (random.randint(0, WIDTH-30), random.randint(0, HEIGHT-30))

pg.draw.rect(screen, (255, 0, 0), item) # Item zeichnen
\end{lstlisting}
\end{myexercise}

\begin{myexercise}[Highscore anzeigen]\label{myexercise:Highscore}
    Erstellen Sie eine Variable \lstinline|score|, die bei jeder Kollision um 1 erhöht wird. Zeichnen Sie den aktuellen Punktestand mit \lstinline|font.render(...)| oben links im Fenster. Verwenden Sie folgenden Code, um Text zu zeichnen:
    \begin{lstlisting}
# Vor der Game-Schleife:
font = pg.font.Font(None, 36)  # Schriftart und -grösse
# ...
# In der Game-Schleife, im Render-Abschnitt:
score_text = font.render(f"Score: {score}", True, (0, 0, 0))
screen.blit(score_text, (10, 10))  # Text oben links zeichnen
\end{lstlisting}
	\begin{myanswer}
		Diese Musterlösung finden Sie in \cref{sec:K09_sol}.
	\end{myanswer}
\end{myexercise}

\begin{myexercise}[Ereignisse (Keyboard / Maus)]\label{myexercise:Ereignisse}
    Reagieren Sie auf \lstinline|KEYDOWN|- und \lstinline|MOUSEBUTTONDOWN|-Events. Beispiel: Bei Mausklick wird an der Klickposition ein kleiner Kreis gezeichnet (oder eine Partikelspur gestartet). Bei \lstinline|ESC| beendet sich das Game. Da die Kreise jedes Frame neu gezeichnet werden müssen, speichern Sie die Positionen in einer Liste.

    Um Ereignisse wie etwa einen Maus-Klick zu verarbeiten, verwenden Sie folgenden Code:
    \begin{lstlisting}
# Vor der Game-Schleife:
mouse_positions = []  # Liste für Maus-Klick-Positionen

# In der Game-Schleife, im Event-Abschnitt:
for event in pg.event.get():
	if event.type == pg.QUIT:
		running = False
	elif event.type == pg.KEYDOWN and event.key == pg.K_ESCAPE:
		running = False
	elif event.type == pg.MOUSEBUTTONDOWN:
		x, y = event.pos
        mouse_positions.append((x, y))  # Position speichern

# In der Game-Schleife, im Render-Abschnitt:
for pos in mouse_positions:
    pg.draw.circle(screen, (255, 100, 100), pos, 10)
\end{lstlisting}
	\begin{myanswer}
		Diese Musterlösung finden Sie in \cref{sec:K09_sol}.
	\end{myanswer}
\end{myexercise}

\begin{myexercise}[Medien: Bild und Sound]\label{myexercise:Medien}
    Laden Sie ein Bild (\lstinline|.png /.jpg|) und zeichnen Sie es mit \lstinline|screen.blit(...)|. Skalieren/rotieren Sie es optional. Laden Sie einen Sound (\lstinline|.wav/ .ogg / .mp3|) mit \lstinline|pg.mixer.Sound| und spielen Sie ihn bei einer Kollision ab.

    Für Bilder:
    \begin{lstlisting}
# Vor der Game-Schleife:
img = pg.image.load("assets/player.jpeg") # Bild laden
img = pg.transform.scale(img, (64, 64)) # optional: skalieren

# In der Game-Schleife, im Render-Abschnitt:
screen.blit(img, player) # Bild an Spieler-Position zeichnen
\end{lstlisting}

    Für Sounds:
    \begin{lstlisting}
# Vor der Game-Schleife:
pg.mixer.init() # Mixer initialisieren (ist nötig für Sound)
motor = pg.mixer.Sound("assets/motor.wav") # Sound laden

# In der Game-Schleife, bei Kollision:
if player.colliderect(item):
    motor.play() # Sound abspielen
\end{lstlisting}

    Hinweis: Legen Sie die Medien in einem Ordner (z.\,B. \lstinline|assets/|) ab und achten Sie auf relative Pfade. Die Dateien \texttt{player.jpeg} und \texttt{motor.wav} finden Sie auf Moodle.

    Optional: Skalieren Sie das Bild auf eine bestimmte Höhe, wobei das Seitenverhältnis beibehalten wird, indem Sie das Seitenverhältnis mit den Befehlen \lstinline|img.get_width()| und \lstinline|img.get_height()| berechnen.

    \textbf{Tipp:} Viele kostenlose Soundeffekte finden Sie unter \url{https://pixabay.com/sound-effects}.
	\begin{myanswer}
		Diese Musterlösung finden Sie in \cref{sec:K09_sol}.
	\end{myanswer}
\end{myexercise}

\begin{mychallenge}[Bonus: Kleines Sammelspiel]
    Verbinden Sie die vorherigen Bausteine: Bewegen Sie den Spieler, sammeln Sie Items (Punkte zählen), spielen Sie dabei einen Sound ab und zeichnen Sie im Hintergrund Ihren Sonnenuntergang. Begrenzen Sie die Spielzeit auf 60 Sekunden und zeigen Sie die verbleibende Zeit im Fenstertitel an.
\end{mychallenge}

\begin{myexercise}[Pong-Spiel]
    Erstellen Sie das klassische Pong-Spiel mit zwei Paddles und einem Ball.

    Das folgende Video zeigt das fertige Pong-Spiel: \url{https://youtu.be/vCoBgJlUg_c}

    Die Paddles werden mit den Tasten \lstinline|W/S| (links) und \lstinline|UP/DOWN| (rechts) gesteuert. Der Ball bewegt sich automatisch und prallt von den Paddles und den oberen/unteren Bildschirmrändern ab. Zählen Sie die Punkte, wenn ein Spieler den Ball am Gegner vorbei spielt. Gehen Sie Schritt für Schritt vor, indem Sie folgende Elemente umsetzen:
    \begin{enumerate}
        \item Rechtecke für Paddles und Ball zeichnen
        \item Mittellinie zeichnen
        \item Tastatur-Eingaben für Paddle-Bewegung (\texttt{W/S} und \texttt{UP/DOWN})
        \item Ball von der Mitte aus in eine zufällige Richtung starten lassen
        \item Kollisionserkennung für Ball und Paddles, Richtung von Ball umkehren bei Kollision oder sofern der Ball die oberen/unteren Ränder berührt
        \item Punkte zählen und anzeigen
        \item Soundeffekte bei Kollisionen und Punkten
    \end{enumerate}
    \begin{myanswer}
        Eine mögliche Lösung zum Pong-Spiel finden Sie auf Moodle.
    \end{myanswer}
\end{myexercise}


\begin{myattention}[Hinweise zu Pong]
    Die folgende Sammlung von Hinweisen zeigt kleine Bausteine, die im Pong-Beispiel verwendet werden und Ihnen beim Verständnis oder bei eigenen Varianten helfen.

    \begin{itemize}

        \item \textbf{Zufällige Start-Richtung für den Ball:}
              \begin{lstlisting}
import random
...
ball_speed = 5
# Geschwindigkeit als separate Variablen
ball_speed_x = random.choice((-1, 1)) * ball_speed
ball_speed_y = random.choice((-1, 1)) * ball_speed

# In der Game-Schleife:
ball.x += ball_speed_x
ball.y += ball_speed_y
\end{lstlisting}

        \item \textbf{Rect-Ränder nutzen und Kollisionen \enquote{sauber} auflösen:}
              \begin{lstlisting}
if ball.colliderect(left_paddle):
    ball_speed_x *= -1     # x-Richtung umkehren (Abprall)
\end{lstlisting}

        \item \textbf{Abprall oben/unten (Wände):}
              \begin{lstlisting}
if ball.top <= 0 or ball.bottom >= HEIGHT:
    ball_speed_y *= -1
\end{lstlisting}

        \item \textbf{Trefferwinkel variieren:} Je weiter oben/unten am Paddle getroffen wird, desto mehr vertikale Geschwindigkeit.
              \begin{lstlisting}
offset = (ball.centery - left_paddle.centery) / (paddle_h / 2)  # -1..+1
ball_speed_y = ball_speed * offset
\end{lstlisting}

        \item \textbf{Ball nach Punkt neu zentrieren (als Funktion):}
              \begin{lstlisting}
def reset_ball(direction):  # direction: -1 nach links, +1 nach rechts
    ball.center = (WIDTH // 2, HEIGHT // 2)
    global ball_speed_x, ball_speed_y # Zugriff auf globale Variablen
    ball_speed_x = direction * ball_speed
    ball_speed_y = random.choice((-1, 1)) * ball_speed
\end{lstlisting}

        \item \textbf{Nützliche Rect-Eigenschaften:} \lstinline|left, right, top, bottom, center, centery| helfen bei Ausrichtung und Kollisionen.
    \end{itemize}
\end{myattention}

\begin{mychallenge}[Pong-Bonus]
    Fügen Sie dem Pong-Spiel folgende Features hinzu:
    \begin{itemize}
        \item Soundeffekte bei Paddle-Kollision und Punktgewinn
        \item Startbildschirm mit Anweisungen
        \item Game-Over-Bildschirm nach einer bestimmten Punktzahl (z.\,B. 10 Punkte)
        \item Paddle-Geschwindigkeit erhöhen, wenn der Ball getroffen wird
        \item Hintergrundmusik während des Spiels
    \end{itemize}
\end{mychallenge}

\begin{myexercise}[Videoaufnahme des Spiels]
    Nehmen Sie ein kurzes Video (ca. 1-2 Minuten) auf, das das Gameplay Ihres Spiels zeigt. Laden Sie das Video auf eine Plattform (z.\,B. YouTube, Vimeo) hoch und fügen Sie den Link auf Moodle ein. Für die Bildschirmaufnahme können Sie folgende Tools verwenden:

    \begin{itemize}
        \item \textbf{Windows:} Verwenden Sie die Bildschirmaufnahme-Funktion von PowerPoint (unter \enquote{Einfügen} $\rightarrow$ \enquote{Bildschirmaufnahme}) oder die Xbox Game Bar (\keys{\faWindows + G}), beenden Sie die Aufnahme mit \keys{\cmd + \faWindows + R}
        \item \textbf{macOS:} Verwenden Sie die integrierte Bildschirmaufnahme (\keys{\shift + \cmd + 5}), beenden Sie die Aufnahme mit \keys{\cmd + \ctrl + Esc}
    \end{itemize}



\end{myexercise}

\section{Game-Auftrag}

\subsection{Thema}
Im Folgenden wählen Sie in Zweier- bis Dreier-Gruppen ein Game-Thema aus und entwickeln ein kleines 2D-Game mit \texttt{pygame-ce}. Mögliche Themen sind:
\begin{itemize}
    \item \textbf{Sammelspiel:} Bewegen Sie eine Spielfigur (z.\,B. Auto, Raumschiff) und sammeln Sie Items ein (z.\,B. Treibstoff, Sterne). Jedes eingesammelte Item gibt Punkte.
    \item \textbf{Endlos-Runner:} Steuern Sie eine Spielfigur (z.\,B. Läufer, Auto) und weichen Sie Hindernissen aus (z.\,B. Mauern, andere Autos). Jedes überstandene Hindernis gibt Punkte.
    \item \textbf{Fangspiel:} Steuern Sie eine Spielfigur (z.\,B. Katze, Netz) und fangen Sie herumspringende Objekte (z.\,B. Mäuse, Schmetterlinge). Jedes gefangene Objekt gibt Punkte.
    \item \textbf{Shooter (z.B. Space Invaders):} Steuern Sie ein Raumschiff und schiessen Sie auf herannahende Gegner. Jeder abgeschossene Gegner gibt Punkte.
\end{itemize}

Weitere Themen können in Absprache mit der Lehrperson gewählt werden. Wichtig ist, dass das Spiel die geforderten Python-Konzepte sinnvoll einsetzt (siehe Anforderungen) und die Komplexität angemessen ist (nicht zu einfach, aber auch nicht zu komplex).

\subsection{Anforderungen}

Ihr Projekt sollte folgende, essentielle Python-Konzepte sinnvoll einsetzen
\begin{todolist}
    \item Verwendung von \texttt{pygame-ce} für Fenster, Game-Schleife, Ereignisse, Zeichnen, Kollisionen und Medien (Bilder, Sound), Steuerung mit Tastatur und/oder Maus
    \item Verwendung physikalischer Formeln für Bewegungen (z.\,B. Geschwindigkeit, Beschleunigung, Schwerkraft)
    \item Verwendung von \textbf{Variablen} zur Speicherung von Spielzuständen (z.\,B. Spielerposition, Punktestand, verbleibende Zeit)
    \item Verwendung von \textbf{Funktionen} zur Strukturierung des Codes (z.\,B. \lstinline|draw_player()|, \lstinline|update_game()|, \lstinline|handle_input()|), mit sowie ohne \lstinline|return|-Wert
    \item Verwendung von \textbf{logischen Ausdrücken} und \textbf{Schleifen} zur Steuerung des Spielablaufs (z.\,B. Kollisionserkennung, Punkte zählen, Spielzeit)
    \item Verwendung von \textbf{Datenstrukturen} (Listen, Dictionaries) zur Verwaltung von mehreren Spielobjekten (z.\,B. Liste von Items, Liste von Hindernissen)
    \item Verwendung von \textbf{Klassen} zur Modellierung von Spielobjekten (z.\,B. \lstinline|class Player|, \lstinline|class Item|, \lstinline|class Obstacle|)
    \item Verwendung mehrerer Python-Unterdateien zur besseren Strukturierung des Codes (z.\,B. \lstinline|game.py|, \lstinline|player.py|, \lstinline|item.py|)
\end{todolist}

Ihr Game muss \textbf{Kommentare} enthalten, die den Code erklären und die Struktur des Programms verdeutlichen. Verwenden Sie sowohl einzeilige Kommentare (mit \lstinline|#|) als auch mehrzeilige Kommentare (mit \texttt{"""..."""}) für Funktionen und Klassen. \textbf{Schauen Sie sich die detaillierten Bewertungskriterien auf Moodle an, um sicherzustellen, dass Ihr Code den Anforderungen entspricht.}

\subsection{Bonus}

\begin{todolist}
    \item Python-spezifisch:
    \begin{todolist}
        \item Verwendung von \texttt{try/except} zur Behandlung von Fehlern (z.\,B. Laden von Medien, Division durch Null)
        \item Verwendung von \texttt{with}-Anweisung zum sicheren Öffnen und Schliessen von Dateien (z.\,B. Highscore speichern)
        \item Wenn Sie ein Punktesystem mit Highscore (gespeichert in einer Datei) implementieren.
    \end{todolist}
    \item Game-spezifisch:
    \begin{todolist}
        % \item Wenn Sie ein Startmenü und/oder ein Game-Over-Bildschirm implementieren.
        % \item Wenn Sie Animationen (z.\,B. bewegte Spielfigur, animierte Hindernisse) einbauen.
        % \item Wenn Sie verschiedene Level oder Schwierigkeitsgrade (z.\,B. schnellere Hindernisse, mehr Items) implementieren.
        \item Wenn Sie viele Power-Ups (z.\,B. temporäre Unverwundbarkeit, Doppelte Punkte) einbauen.
        \item Wenn Sie viele eigens erstellte Hintergrundmusik und/oder Soundeffekte (z.\,B. bei Kollision, Item-Sammlung) einbauen.
        \item Wenn Sie das Spiel mit einem Gamepad/Controller steuern können.
        \item Wenn Sie einen Online-Multiplayer-Modus (z.\,B. über LAN) implementieren.
    \end{todolist}
    \item Weitere Ideen für Bonus:
    \begin{todolist}
        \item Wenn Sie das Spiel mit einem Level-Editor (zum Erstellen eigener Level) erweitern können.
        \item Wenn Sie das Spiel mit einem Story-Modus (mit Handlung und Charakteren) erweitern können.
        \item Wenn Sie das Spiel mit einem besonders raffinierten Achievements-System (Erfolge) erweitern können.
        \item Wenn Sie das Spiel mit einem Debug-Modus (zum Testen und Entwickeln) erweitern können.
    \end{todolist}
\end{todolist}

% \subsection{Abgabe der Games}
% Im Rahmen dieses Projekts geben Sie drei Berichte ab, die den Fortschritt und die Reflexion Ihres Spiels dokumentieren. Jeder Bericht sollte klar strukturiert und gut lesbar sein. Verwenden Sie Überschriften, Absätze und Aufzählungen, um den Text zu gliedern. Achten Sie auf eine korrekte Rechtschreibung und Grammatik. Die Berichte sollten je als eine einzige PDF-Datei eingereicht werden.
% \begin{enumerate}
%     \item \textbf{Projekt-Entwurf (nach 2 Wochen):} Konzept und Design des Spiels, Herausforderungen, Planung (ca. 1.5--2 Seiten). In diesem Bericht sollte Folgendes enthalten sein:
%           \begin{todolist}
%               \item \textbf{Thema}: Kurzbeschreibung des Spiels (Genre, Setting, Ziel), ca. 4--5 Sätze
%               \item \textbf{Spielmechanik}: Beschreibung der Hauptspielmechanik (Steuerung, Punktesystem, Level), ca. 5--6 Sätze
%               \item \textbf{Medien}: Liste der benötigten Medien (Bilder, Sounds) und deren Quellen (selbst erstellt, frei verfügbar)
%               \item \textbf{Komplexität}: Einschätzung der Komplexität des Spiels (Anzahl der Funktionen, Klassen, Datenstrukturen)
%               \item \textbf{Zeitplan}: Grober Zeitplan für die Umsetzung (Was wird in den nächsten Wochen gemacht?)
%               \item \textbf{Gruppen-Aufteilung}: Wer macht was in der Gruppe?
%               \item \textbf{Herausforderungen}: Welche Herausforderungen erwarten Sie bei der Umsetzung?
%           \end{todolist}
%     \item \textbf{Zwischenbericht (nach 4 Wochen):} Fortschritt und Herausforderungen (ca. 1 Seite). In diesem Bericht sollte Folgendes enthalten sein:
%           \begin{todolist}
%               \item \textbf{Fortschritt}: Beschreibung des aktuellen Stands des Spiels (Was ist bereits umgesetzt? Was fehlt noch?)
%               \item \textbf{Herausforderungen}: Welche Herausforderungen sind aufgetreten und wie wurden sie gelöst?
%               \item \textbf{Zeitplan}: Aktualisierung des Zeitplans für die verbleibende Zeit (Was ist bereits umgesetzt? Was fehlt noch? Welche Probleme sind aufgetreten?)
%               \item \textbf{Gruppen-Aufteilung}: Hat sich die Aufgabenverteilung in der Gruppe geändert? Wie wird die kommende Arbeit aufgeteilt?
%           \end{todolist}
%     \item \textbf{Endbericht (nach 6 Wochen)}: Fertigstellung und Reflexion des Spiels (ca. 2--3 Seiten). In diesem Abschlussbericht sollte Folgendes enthalten sein:
%           \begin{todolist}
%               \item \textbf{Spielbeschreibung}: Kurzbeschreibung des fertigen Spiels (Genre, Setting, Ziel), ca. 4--5 Sätze
%               \item \textbf{Spielmechanik}: Beschreibung der Hauptspielmechanik (Steuerung, Punktesystem, Level), ca. 5--6 Sätze
%               \item \textbf{Reflexion}: Was lief gut? Was lief schlecht? Was würden Sie beim nächsten Mal anders machen?
%               \item \textbf{Lernziele}: Welche neuen Konzepte/Techniken haben Sie gelernt und angewendet?
%               \item \textbf{Gruppenarbeit}: Wie hat die Zusammenarbeit in der Gruppe funktioniert? \textbf{Weisen Sie klar aus, wer für welche Teile des Codes verantwortlich ist.} Jede Zeile Code muss mindestens einer Person zugeordnet sein!
%               \item \textbf{Herausforderungen}: Welche Herausforderungen sind aufgetreten und wie wurden sie gelöst?
%               \item \textbf{Medien}: Liste der verwendeten Medien (Bilder, Sounds) und deren Quellen
%               \item \textbf{Bonus}: Welche Bonuspunkte wurden erreicht?
%               \item \textbf{Screenshots}: Fügen Sie 2--3 Screenshots des Spiels bei, die verschiedene Spielsituationen zeigen.
%               \item \textbf{Screen Recording}: Nehmen Sie ein kurzes Video (ca. 1--2 Minuten) auf, das das Gameplay zeigt. Laden Sie das Video auf eine Plattform (z.B. YouTube, Vimeo) hoch und fügen Sie den Link im Bericht ein.
%               \item \textbf{Quellcode}: Fügen Sie den Link zum GitHub-Repository bei, in dem der Quellcode des Spiels gespeichert ist.
%           \end{todolist}
% \end{enumerate}

% Die endgültige Abgabe des Games erfolgt via GitHub. Erstellen Sie ein öffentliches Repository mit dem Quellcode, den Medien und einer \texttt{README.md}-Datei, die das Spiel beschreibt und Anweisungen zur Ausführung enthält. Reichen Sie den Link zum Repository via Moodle ein. Für Anleitungen zur Erstellung eines GitHub-Repositories und zum Hochladen von Code, s. Anhang \ref{ch:GitHub}.

\section{Bewertung}
\subsection{Projektbewertung}
Die Bewertung Ihres Projekts erfolgt anhand der auf Moodle aufgelisteten Kriterien. Die maximale Punktzahl beträgt 100 Punkte.
% \begin{itemize}
%     \item \textbf{Funktionalität (25\%):} Das Spiel läuft fehlerfrei und erfüllt die grundlegenden Anforderungen (Steuerung, Kollisionen, Punktesystem).
%     \item \textbf{Code-Qualität (25\%):} Der Quellcode ist gut strukturiert, lesbar und verwendet die geforderten Python-Konzepte (Variablen, Funktionen, Klassen, Datenstrukturen etc.).
%     \item \textbf{Kreativität (20\%):} Das Spiel zeigt Originalität und Kreativität in Design und Umsetzung (Thema, Spielmechanik, Medien).
%     \item \textbf{Dokumentation (30\%):} Die Berichte sind klar, gut strukturiert und reflektieren den Entwicklungsprozess.
% \end{itemize}

\subsection{Gruppen-Besprechung des Spiels}

Im Anschluss an Ihre Abgabe findet eine Besprechung Ihres Spiels in Gruppen mit der Lehrperson statt (ca. 15 Minuten pro Gruppe). Dabei werden Ihnen Fragen zu Ihrem Spiel gestellt, um Ihr Verständnis und Ihre Reflexion zu überprüfen. Achten Sie darauf, in Ihrem Schlussbereicht zu erwähnen, wer für welche Teile zuständig war -- das Verständnis von jedem Gruppenmitglied wird speziell überprüft. Die Besprechung resultiert in einer \textbf{individuellen} Verständnis-Prozentzahl von 0\%--100\% Punkten, welche mit Ihrer Endnote multipliziert wird.

Ihre Endnote wird wie folgt berechnet:
\[\text{Endnote} = 5 \cdot \lr{ \text{Projekt-Punkte} \cdot \lr{\frac{\text{Verständnis-Prozentzahl}}{100}}} + 1\]

\begin{myattention}[KI-Tools]
    Es wird von Ihnen erwartet, dass Sie sich aktiv an der Entwicklung des Spiels beteiligen und die geforderten Python-Konzepte verstehen und anwenden können. ChatGPT oder andere KI-Tools dürfen Sie verwenden, solange Sie Ihren gesamten Code verstehen und erklären können. Wie Sie im obigen Benotungsschema erkennen können, kann Ihre Note durch mangelndes Verständnis stark beeinträchtigt werden (bis zur Note 1). Falls Sie ChatGPT oder andere KI-Tools verwenden, müssen Sie dies im Code klar dokumentieren (z.B. mit Kommentaren oder in einer Begleit-Dokumentation). Jede Zeile Code muss mindestens einer Person zugeordnet sein!
\end{myattention}


\iftoggle{exerciseonly}{}{%
	\clearpage

	\section{Längere Lösungen}\label{sec:K09_sol}

	\begin{myanswer}[myexercise:Sonnenuntergang]
		\leavevmode %
        \lstinputlisting[language=Python,caption=\texttt{game\_intro\_sunset.py},label=listing:Sonnenuntergang]{Grundlagen_Info/00_Programmieren/Skript/Code/K08_Game/game_intro_sunset.py}
	\end{myanswer}

    \begin{myanswer}[myexercise:Bewegungen]
		\leavevmode %
        \lstinputlisting[language=Python,caption=\texttt{game\_intro\_player\_rect.py},label=listing:Bewegungen]{Grundlagen_Info/00_Programmieren/Skript/Code/K08_Game/game_intro_player_rect.py}
	\end{myanswer}

    \begin{myanswer}[myexercise:Kollisionen]
        \leavevmode %
        \lstinputlisting[language=Python,caption=\texttt{game\_intro\_collision.py},label=listing:Kollisionen]{Grundlagen_Info/00_Programmieren/Skript/Code/K08_Game/game_intro_collision.py}
    \end{myanswer}

    % \begin{myanswer}[myexercise:Highscore]
    %     \leavevmode %
    %     \lstinputlisting[language=Python,caption=\texttt{TODO},label=listing:Highscore]{Grundlagen_Info/00_Programmieren/Skript/Code/K08_Game/TODO}
    % \end{myanswer}

    \begin{myanswer}[myexercise:Ereignisse]
        \leavevmode %
        \lstinputlisting[language=Python,caption=\texttt{game\_intro\_events.py},label=listing:Ereignisse]{Grundlagen_Info/00_Programmieren/Skript/Code/K08_Game/game_intro_events.py}
    \end{myanswer}

    \begin{myanswer}[myexercise:Medien]
        \leavevmode %
        \lstinputlisting[language=Python,caption=\texttt{game\_intro\_media.py},label=listing:Medien]{Grundlagen_Info/00_Programmieren/Skript/Code/K08_Game/game_intro_media.py}
    \end{myanswer}
}